\section{Discussion}\label{sec-discussion}

\subsection{Le verrou visé est bien celui qui est levé}

La lecture centrale du \cref{tab-comparaison} est le gradient des gains
selon la taille des objets : $+469{,}4\,\%$ de rappel pour les objets
petits, $+135{,}6\,\%$ pour les moyens, $+20{,}2\,\%$ pour les grands.
L'amélioration est d'autant plus forte que l'objet est petit, c'est-à-dire
exactement là où le modèle générique échouait : le modèle pré-entraîné ne
retrouvait qu'un objet petit sur huit ($0{,}122$), le modèle fine-tuné en
retrouve plus des deux tiers ($0{,}694$). La réponse à la question de
recherche est donc positive, et elle l'est de manière ciblée : le
fine-tuning sur du trafic urbain dense n'améliore pas seulement les
métriques globales, il lève spécifiquement le verrou des véhicules
éloignés, de petite taille ou masqués qui motivait ce travail. Les
augmentations orientées petits objets (\emph{mosaic}, \emph{scale},
\cref{sec-entrainement}) et la densité du jeu d'entraînement
($8{,}5$ véhicules par image en moyenne) sont les explications les plus
plausibles de ce ciblage, sans que notre protocole permette d'isoler la
contribution de chacune ; une étude d'ablation serait nécessaire.

\subsection{Cohérence avec l'écart de domaine rapporté par BMD-45}

Notre passage de $0{,}296$ à $0{,}644$ de mAP@0,5:0,95, soit un facteur
$2{,}2$, est cohérent avec l'écart de domaine rapporté par les auteurs de
BMD-45~\cite{sharma2026bmd45} : des détecteurs entraînés sur
UA-DETRAC~\cite{wen2020uadetrac} atteignent $33{,}6\,\%$ de mAP@0,5:0,95
sur leur jeu, contre $83{,}8\,\%$ pour des modèles entraînés en domaine,
soit un facteur $2{,}5$. Nous retrouvons cet ordre de grandeur avec un
modèle nano et $6{,}7\,\%$ des images du jeu complet : l'essentiel de
l'écart entre un détecteur générique et un détecteur adapté semble donc
se combler dès un fine-tuning modeste, ce qui est une information
directement utile pour un déploiement à budget contraint.

\subsection{Ce que le fine-tuning ne résout pas}

Le rappel des objets petits plafonne à $0{,}694$ : 246 des 804 objets
petits du jeu de test restent non détectés. Une partie de ce résidu
relève d'une limite informationnelle et non d'un défaut de l'approche :
sous un certain nombre de pixels, l'information nécessaire à la détection
n'existe simplement plus dans l'image, et aucun modèle, quelle que soit
sa taille, ne peut la reconstituer. Ce plafond suggère que les marges de
progression restantes sont à chercher du côté de la résolution d'entrée,
de l'inférence par tuiles ou du placement des caméras, plutôt que du seul
entraînement.

Le gain de débit ($+16{,}2\,\%$, de $37{,}2$ à $43{,}2$ FPS) doit être
interprété avec prudence : l'architecture étant identique, il ne peut
provenir du réseau lui-même. L'hypothèse la plus plausible est un
post-traitement (suppression des non-maxima) allégé par la réduction de
80 à 4 classes, à quoi s'ajoute la variabilité de mesure inhérente à un
environnement partagé comme Google Colab. Nous n'avons pas conçu
d'expérience pour départager ces deux effets.

\subsection{Limites}

\paragraph{Domaine indien, contexte béninois.}
Il n'existe à ce jour aucun jeu de données public de vision par
ordinateur spécifique au Bénin. Nous avons donc travaillé sur du trafic
dense indien, structurellement comparable au trafic béninois sur trois
points : caméras CCTV fixes, forte densité, dominance des deux-roues
($56\,\%$ des instances, à rapprocher des zémidjans). Nos résultats
établissent que l'adaptation à un domaine de trafic dense de ce type est
possible et fortement bénéfique ; ils ne démontrent pas que le modèle
obtenu est adapté au trafic béninois lui-même. La constitution d'un jeu
de données local reste le préalable à toute conclusion de ce type.

\paragraph{Échelle des expériences.}
Le sous-ensemble utilisé représente 3\,000 images sur les 45\,000
disponibles, le modèle est la déclinaison nano, et l'entraînement est
limité à 50 epochs : ces choix, imposés par le temps machine disponible,
bornent la portée des chiffres absolus rapportés, même si la comparaison
relative entre les deux modèles reste valide puisqu'ils sont évalués dans
des conditions identiques.

\paragraph{Classes minoritaires.}
Les bus (189 instances) et camions (240 instances) sont minoritaires
dans le jeu de test : les métriques qui les concernent sont
mécaniquement moins stables que celles des voitures (742) et motos
(1\,380), et nous ne rapportons pas de métriques par classe pour cette
raison.

\paragraph{Contamination potentielle entre splits.}
Nos splits sont dérivés de façon déterministe, mais BMD-45 étant
constitué d'images de caméras fixes, des images quasi identiques
(même caméra, instants proches) peuvent se retrouver de part et d'autre
de la séparation train/test, un risque de contamination connu dans les
jeux publics qui tendrait à flatter les métriques absolues des deux
modèles sans invalider leur comparaison.

\paragraph{Protection des données personnelles.}
Les plaques d'immatriculation et les visages sont des données à
caractère personnel. Tout déploiement réel d'un système de
vidéoprotection au Bénin suppose un cadre juridique conforme au code du
numérique~\cite{benin2017code} et le contrôle de l'Autorité de
protection des données personnelles (APDP). Le présent travail, mené sur
un jeu de données public diffusé sous licence CC~BY~4.0, n'implique
aucune collecte de données.
